% XVI Portuguese Category Seminar | Coimbra, 12--14 November 2026
% Compile with pdfLaTeX (for example in Overleaf or TeX Live).
% This file is self-contained: no external images or bibliography file.
% Replace the example metadata, abstract and references below.
% Email Maria Manuel Clementino at mmc@mat.uc.pt.
% Subject: PCS-2026 SUBMISSION
% Paste the completed registration form from the website into your email
% body and attach this .tex file and its compiled PDF.
% Abstract deadline: 29 October 2026, 23:59 Coimbra time.
% Attendance-only registration: 5 November 2026, 23:59 Coimbra time.
% Decisions will be communicated by 3 November 2026.

\documentclass[11pt,a4paper]{article}
\usepackage[T1]{fontenc}
\usepackage[utf8]{inputenc}
\usepackage{lmodern}
\usepackage[margin=27mm]{geometry}
\usepackage{amsmath,amssymb,amsthm}
\usepackage{microtype}
\usepackage{xcolor}
\usepackage{enumitem}
\usepackage[colorlinks=true,urlcolor=PCSGold,citecolor=PCSGold,
            linkcolor=PCSGold,unicode=true]{hyperref}
\definecolor{PCSGold}{HTML}{926009}
\definecolor{PCSGray}{HTML}{565A53}
\setlength{\parindent}{0pt}
\setlength{\parskip}{0.65em}
\setlength{\emergencystretch}{2em}
\pagestyle{empty}

% ---------------------- YOUR INFORMATION ----------------------
\newcommand{\talktitle}{Title of the proposed talk}
\newcommand{\talkauthors}{First Author\textsuperscript{1}}
\newcommand{\talkaffiliations}{%
  \textsuperscript{1}Department, University or Institution, Country}
\newcommand{\talkemail}{author@example.org}
\newcommand{\talkkeywords}{Category theory; keyword two; keyword three}
\newcommand{\talksubject}{Theory of Categorical Structures}
\newcommand{\talkmsc}{18A05} % Optional: leave empty to omit this line.

% Multiple authors: replace the definitions above, for example:
% \newcommand{\talkauthors}{First Author\textsuperscript{1},
%                           Second Author\textsuperscript{2}}
% \newcommand{\talkaffiliations}{%
%   \textsuperscript{1}First Institution, Country\\
%   \textsuperscript{2}Second Institution, Country}
% Give the presenting author's email in \talkemail.
% Choose the closest subject area from the registration form.
% MSC2020 is optional; list the primary code first, then secondary codes.

\hypersetup{pdftitle={\talktitle},pdfsubject={XVI Portuguese Category Seminar}}

\begin{document}
{\small\sffamily\color{PCSGold}
 XVI PORTUGUESE CATEGORY SEMINAR\hfill COIMBRA 2026\par}
\vspace{0.2em}
{\small\sffamily\color{PCSGray}12--14 November 2026\par}
{\color{PCSGold}\rule{\linewidth}{0.7pt}}

\vspace{1.5em}
{\LARGE\bfseries \talktitle\par}
\vspace{0.6em}
{\large \talkauthors\par}
\vspace{0.2em}
{\small\color{PCSGray}\talkaffiliations\par}
{\small\href{mailto:\talkemail}{\talkemail}\par}

\vspace{1.1em}
{\large\bfseries Abstract\par}

% ------------------------- YOUR ABSTRACT ----------------------
Replace this paragraph with your abstract. Introduce the problem and its
context, state the main result or question, and describe the contribution
of the proposed talk. The text should be understandable to researchers
working in neighbouring areas of category theory. Replace or remove all
example text, keywords, classifications and references before submission.

You may use standard mathematical notation. For example, an adjunction
between categories $\mathcal{C}$ and $\mathcal{D}$ consists of functors
$F\colon\mathcal{C}\to\mathcal{D}$ and
$G\colon\mathcal{D}\to\mathcal{C}$ together with a natural isomorphism
\[
  \operatorname{Hom}_{\mathcal{D}}(FX,Y)
  \cong
  \operatorname{Hom}_{\mathcal{C}}(X,GY).
\]
This sentence illustrates a bibliographic citation~\cite{maclane}.
Include acknowledgements of joint work or funding here if appropriate.

\vspace{0.7em}
{\small
 \textbf{Keywords:} \talkkeywords\par
 \textbf{Subject area:} \talksubject\par
 \ifx\talkmsc\empty\else
   \textbf{MSC2020:} \talkmsc\par
 \fi
}

% ------------------------- BIBLIOGRAPHY ------------------------
% Use \bibitem entries below; cite them in the abstract with \cite{key}.
% Remove the example reference if it is not relevant to your abstract.
\renewcommand{\refname}{References}
\begin{thebibliography}{9}
\small
\bibitem{maclane}
S. Mac Lane, \emph{Categories for the Working Mathematician},
2nd edition, Graduate Texts in Mathematics \textbf{5}, Springer, 1998.
\end{thebibliography}

\end{document}
